\documentclass[12pt,a4paper]{article}

% ================= PACOTES =================
\usepackage[utf8]{inputenc}
\usepackage[T1]{fontenc}
\usepackage{amsmath,amssymb,amsfonts}
\usepackage{geometry}
\usepackage{graphicx}
\usepackage{float}
\usepackage{hyperref}
\usepackage{xcolor}
\usepackage{booktabs}
\usepackage{enumitem}
\usepackage{fancyhdr}
\usepackage{titlesec}
\usepackage{caption}

\geometry{margin=2.5cm}

\hypersetup{
    colorlinks=true,
    linkcolor=blue,
    citecolor=blue,
    urlcolor=blue
}

\pagestyle{fancy}
\fancyhf{}
\rhead{APC \#2 -- Cálculo Numérico Aplicado}
\lhead{\thepage}
\renewcommand{\headrulewidth}{0.4pt}

\titleformat{\section}{\normalfont\Large\bfseries}{\thesection}{1em}{}
\titleformat{\subsection}{\normalfont\large\bfseries}{\thesubsection}{1em}{}

% ================= DADOS DO ALUNO =================

\title{
    \vspace{-1.5cm}
    \textbf{Atividade para Casa 02} \\}

\author{Elisa Vargas de Souza \\ Matrícula: 231012540 \\[0.3cm]
    Cálculo Numérico Aplicado \\}

\begin{document}

\maketitle

% =====================================================
\section{Polinômio Característico}
% =====================================================

Considerando o sistema massa-mola-amortecedor com N graus de liberdade, temos que sua equação do movimento é:

\begin{equation}
    M\ddot{\mathbf{x}}(t) + C\dot{\mathbf{x}}(t) + K\mathbf{x}(t) = \mathbf{0}
    \label{eq:eom}
\end{equation}

\noindent em que $M, C, K$ são, respectivamente, as
matrizes de massa, amortecimento e rigidez, e $\mathbf{x}(t)$ é o vetor de deslocamentos. Assumimos então a solução $\mathbf{x}(t) = \boldsymbol{\phi}\, e^{\lambda t}$ e, substituindo em (1), temos:

\begin{equation}
    M\frac{d^2}{dt^2}(\boldsymbol{\phi}\, e^{\lambda t}) + C\frac{d}{dt}(\boldsymbol{\phi}\, e^{\lambda t}) + K(\boldsymbol{\phi}\, e^{\lambda t}) = 0
\end{equation}
\begin{equation}
    \frac{d}{dt}(\boldsymbol{\phi}\, e^{\lambda t}) = \boldsymbol{\phi}\lambda\, e^{\lambda t}, \qquad
    \frac{d^2}{dt^2}(\boldsymbol{\phi}\, e^{\lambda t}) = \boldsymbol{\phi}\lambda^2\, e^{\lambda t}
\end{equation}
\begin{equation}
    M(\boldsymbol{\phi}\lambda^2\, e^{\lambda t}) + C(\boldsymbol{\phi}\lambda\, e^{\lambda t}) + K(\boldsymbol{\phi}\, e^{\lambda t}) = \mathbf{0}
    \label{eq:eom}
\end{equation}

\noindent Sabendo que $\, e^{\lambda t}\neq0$ temos:

\begin{equation}
    M\phi\lambda^2 + C\phi\lambda + K\phi = \mathbf{0}
    \label{eq:eom}
\end{equation}

\noindent Esta equação, no entanto, é a forma geral da operação matricial do problemas. Definindo essa equação como uma única matriz temos:

\begin{equation}
    \phi A(\lambda) = M\phi\lambda^2 + C\phi\lambda + K\phi
    \label{eq:eom}
\end{equation}

\noindent Em que:

\begin{equation}
    A(\lambda) = 
    \begin{bmatrix}
        m_{11}\lambda^2+c_{11}\lambda+k_{11} & m_{12}\lambda^2+c_{12}\lambda+k_{12} & ... \\
        m_{21}\lambda^2+c_{21}\lambda+k_{21} & ... & ...\\
        ... & ... & ...\\
        ... & ... & m_{NN}\lambda^2+c_{NN}\lambda+k_{NN}
    \end{bmatrix}
    \label{eq:eom}
\end{equation}

\noindent Sendo assim, para que hajam soluções para esse sistema que não sejam $A(\lambda) = 0$, o determinante da matriz deve ser igual a 0, ou seja:

\begin{equation}
    det(A(\lambda)) = (m_{11}\lambda^2+c_{11}\lambda+k_{11})(m_{22}\lambda^2+c_{22}\lambda+k_{22})...(m_{NN}\lambda^2+c_{NN}\lambda+k_{NN})
\end{equation}
\begin{equation}    
    - (m_{1N}\lambda^2+c_{1N}\lambda+k_{1N})...(m_{N1}\lambda^2+c_{N1}\lambda+k_{N1})
\end{equation}
\begin{equation}    
    det(A(\lambda)) = a_{2N}\lambda^{2N} + a_{2N-1}\lambda^{2N-1}+...+a_0 = 0
\end{equation}

\noindent Sendo esse o polinômio característico do problema de grau 2N.

% =====================================================
\section{Interpretação Física}
% =====================================================
% TODO: discuta o significado físico das 2N raízes de P(lambda):
% - raízes complexas conjugadas -> modos oscilatórios amortecidos
% - parte real negativa -> estabilidade / decaimento
% - parte real positiva -> instabilidade
% - raízes reais puras -> comportamento superamortecido / não oscilatório
% - relação entre |Im(lambda)| e a frequência natural amortecida
% - relação entre Re(lambda) e a taxa de decaimento / fator de amortecimento

...

% =====================================================
\section{Determinação dos valores}
% =====================================================

% TODO: escolha N (sugestão: N = 2 ou N = 3) e proponha valores numéricos
% para massas m_i, coeficientes de amortecimento c_i e rigidezes k_i.
% Desenhe (ou descreva) o sistema massa-mola-amortecedor correspondente.

Considerando um sistema com $N = 2$ graus de liberdade, foram escolhidos valores para as matrizes de massa, amortecimento e rigidez com base no modelo de um quarto de carro, muito utilizado para análises de dinâmica veícular. Sendo assim, temos:

\begin{equation}
    m_1 = 350, \quad m_2 = 45
\end{equation}

\noindent Onde $m_1$ representa a massa suspensa e $m_2$ a massa não suspensa do veículo. Conectadas pelas rigidezes:

\begin{equation}
    k_1 = 20000, \quad k_2 = 200000
\end{equation}

\noindent Em que $k_1$ representa a rigidez da suspensão e $k_2$ a rigidez do pneu. E, por fim, com coeficientes de amortecimento:

\begin{equation}
    c_1 = 1500, \quad c_2 = 0
\end{equation}

\noindent Onde $c_1$ representa o amortecimento da suspensão e $k_2$ o amortecimento do pneu, que nesse caso é desprezível.

% Sugestão: inclua aqui uma figura esquemática do sistema (TikZ) se desejar.
% \begin{figure}[H]
%     \centering
%     % \input{seu_diagrama_tikz.tex}
%     \caption{Esquema do sistema massa-mola-amortecedor com $N$ graus de liberdade.}
%     \label{fig:sistema}
% \end{figure}

% =====================================================
\section{Sistema de Equações e Polinômio Caracterísco}
% =====================================================

% TODO: monte explicitamente as matrizes M, C, K para os valores escolhidos

As matrizes de massa, amortecimento e rigidez para o sistema proposto são:

\begin{equation}
    M =
    \begin{bmatrix}
        350 & 0 \\
        0 & 45
    \end{bmatrix}, \qquad
    C =
    \begin{bmatrix}
        1500 & -1500 \\
        -1500 & 1500
    \end{bmatrix}, \qquad
    K =
    \begin{bmatrix}
        20000 & -20000 \\
        -20000 & 220000
    \end{bmatrix}.
    \label{eq:matrizes}
\end{equation}

\end{document}
